\subsection{子群与商群的极大环面}

命题 1. 设 G 与 G' 是两个连通紧李群。

\begin{enumerate}
\def\labelenumi{\alph{enumi})}
\item
  设 \(f: \mathbf{G} \to \mathbf{G}'\) 是李群满态射。\(\mathbf{G}'\) 的极大环面是在 \(f\) 下 \(\mathbf{G}\) 的诸极大环面的像。若 \(f\) 的核在 \(\mathbf{G}\) 中是中心的（例如是离散的），则 \(\mathbf{G}\) 的极大环面是在 \(f\) 下 \(\mathbf{G}'\) 的诸极大环面的逆像。
\item
  设 \(\mathrm{H}\) 是 \(\mathrm{G}\) 的连通闭子群。\(\mathrm{H}\) 的每个极大环面都是 \(\mathrm{H}\) 与 \(\mathrm{G}\) 的某个极大环面之交。
\item
  设 H 是 G 的连通闭正规子群。H 的极大环面是 G 的极大环面与 H 的交。
\item
  设 T 是 G 的极大环面；则 L(T) 是 L(G) 的嘉当子代数（第 2 节，定理 2 a)），于是 L(f(T)) 是 L(G') 的嘉当子代数（第 VII 章，§2，第 1 节，命题 4 的推论 2）；由此得出 f(T) 是 G' 的极大环面（第 2 节，定理 2 a)）。若 Ker f 含于 G 的中心，则它含于 T（定理 2 的推论 2），故 T = \(f^{-1}(f(T))\) 。
\end{enumerate}

反之，设 \(T'\) 是 \(G'\) 的一个极大环面；我们证明存在 G 的极大环面 T 使得 \(f(\mathrm{T}) = \mathrm{T}'\) 。设 \(T_{1}\) 是 G 的一个极大环面；则 \(f(\mathrm{T}_{1})\) 是 \(G'\) 的极大环面，且存在 \(g' \in G'\) 使得 \(\mathrm{T}' = g' f(\mathrm{T}_{1}) g'^{-1} (\mathrm{Th.~2~b})\) ；若 \(g \in G\) 满足 \(f(g) = g'\) ，则 \(\mathrm{T}' = f(\mathrm{T})\) ，其中 \(\mathrm{T} = g\mathrm{T}_{1}g^{-1}\) 。

\begin{enumerate}
\def\labelenumi{\alph{enumi})}
\setcounter{enumi}{1}
\item
  设 S 是 H 的一个极大环面；它是 G 的环面，故存在包含 S 的 G 的极大环面 T。于是 \(T \cap H\) 是 H 的包含 S 的交换子群，故等于 S（第 2 节，注 2）。
\item
  由 §1 第 3 节命题 2 c)，L(G) 是 L(H) 与一个理想的直积；于是 L(H) 的嘉当子代数是 L(G) 的嘉当子代数与 L(H) 的交。从而对 G 的任一极大环面 T，T∩H 包含 H 的某个极大环面 S，且 S = T∩H（第 2 节，注 2）。
\end{enumerate}

注。1) 命题 1 立即推广到李代数为紧李代数的连通群。特别地，若 G 是李代数为紧李代数的连通李群，则 G 的嘉当子群（参见第 2 节注 3）恰是连通紧李群 Ad(G) 的诸极大环面（在从 G 到 Ad(G) 的典则同态下）的逆像。

\begin{enumerate}
\def\labelenumi{\arabic{enumi})}
\setcounter{enumi}{1}
\tightlist
\item
  设 G 是连通紧李群，\(\widetilde{\mathrm{D}}(\mathrm{G})\) 是群 \(\mathrm{D}(\mathrm{G})\) 的泛覆盖，\(f : \widetilde{\mathrm{D}}(\mathrm{G}) \to \mathrm{G}\) 是从 \(\tilde{\mathrm{D}}(\mathrm{G})\) 到 \(\mathrm{D}(\mathrm{G})\) 与从 \(\mathrm{D}(\mathrm{G})\) 到 G 的典则态射的复合。则映射 \(T \mapsto f^{-1}(\mathrm{T})\) 是从 G 的极大环面集到 \(\tilde{\mathrm{D}}(\mathrm{G})\) 的极大环面集上的双射；逆双射把极大环面 \(\tilde{T}\) （\(\tilde{\mathrm{D}}(\mathrm{G})\) 中的）对应为 G 的极大环面 \(\mathrm{C}(\mathrm{G})_{0}.f(\tilde{\mathrm{T}})\) 。
\end{enumerate}

